\documentclass[]{article}
\usepackage{lmodern}
\usepackage{amssymb,amsmath}
\usepackage{ifxetex,ifluatex}
\usepackage{fixltx2e} % provides \textsubscript
\ifnum 0\ifxetex 1\fi\ifluatex 1\fi=0 % if pdftex
  \usepackage[T1]{fontenc}
  \usepackage[utf8]{inputenc}
\else % if luatex or xelatex
  \ifxetex
    \usepackage{mathspec}
  \else
    \usepackage{fontspec}
  \fi
  \defaultfontfeatures{Ligatures=TeX,Scale=MatchLowercase}
\fi
% use upquote if available, for straight quotes in verbatim environments
\IfFileExists{upquote.sty}{\usepackage{upquote}}{}
% use microtype if available
\IfFileExists{microtype.sty}{%
\usepackage[]{microtype}
\UseMicrotypeSet[protrusion]{basicmath} % disable protrusion for tt fonts
}{}
\PassOptionsToPackage{hyphens}{url} % url is loaded by hyperref
\usepackage[unicode=true]{hyperref}
\PassOptionsToPackage{usenames,dvipsnames}{color} % color is loaded by hyperref
\hypersetup{
            pdftitle={Overview of the NetSim package},
            pdfauthor={Gustavo A. Ballen},
            colorlinks=true,
            linkcolor=Maroon,
            citecolor=Blue,
            urlcolor=blue,
            breaklinks=true}
\urlstyle{same}  % don't use monospace font for urls
\usepackage[margin=1in]{geometry}
\usepackage{natbib}
\bibliographystyle{apalike}
\usepackage{color}
\usepackage{fancyvrb}
\newcommand{\VerbBar}{|}
\newcommand{\VERB}{\Verb[commandchars=\\\{\}]}
\DefineVerbatimEnvironment{Highlighting}{Verbatim}{commandchars=\\\{\}}
% Add ',fontsize=\small' for more characters per line
\usepackage{framed}
\definecolor{shadecolor}{RGB}{248,248,248}
\newenvironment{Shaded}{\begin{snugshade}}{\end{snugshade}}
\newcommand{\KeywordTok}[1]{\textcolor[rgb]{0.13,0.29,0.53}{\textbf{#1}}}
\newcommand{\DataTypeTok}[1]{\textcolor[rgb]{0.13,0.29,0.53}{#1}}
\newcommand{\DecValTok}[1]{\textcolor[rgb]{0.00,0.00,0.81}{#1}}
\newcommand{\BaseNTok}[1]{\textcolor[rgb]{0.00,0.00,0.81}{#1}}
\newcommand{\FloatTok}[1]{\textcolor[rgb]{0.00,0.00,0.81}{#1}}
\newcommand{\ConstantTok}[1]{\textcolor[rgb]{0.00,0.00,0.00}{#1}}
\newcommand{\CharTok}[1]{\textcolor[rgb]{0.31,0.60,0.02}{#1}}
\newcommand{\SpecialCharTok}[1]{\textcolor[rgb]{0.00,0.00,0.00}{#1}}
\newcommand{\StringTok}[1]{\textcolor[rgb]{0.31,0.60,0.02}{#1}}
\newcommand{\VerbatimStringTok}[1]{\textcolor[rgb]{0.31,0.60,0.02}{#1}}
\newcommand{\SpecialStringTok}[1]{\textcolor[rgb]{0.31,0.60,0.02}{#1}}
\newcommand{\ImportTok}[1]{#1}
\newcommand{\CommentTok}[1]{\textcolor[rgb]{0.56,0.35,0.01}{\textit{#1}}}
\newcommand{\DocumentationTok}[1]{\textcolor[rgb]{0.56,0.35,0.01}{\textbf{\textit{#1}}}}
\newcommand{\AnnotationTok}[1]{\textcolor[rgb]{0.56,0.35,0.01}{\textbf{\textit{#1}}}}
\newcommand{\CommentVarTok}[1]{\textcolor[rgb]{0.56,0.35,0.01}{\textbf{\textit{#1}}}}
\newcommand{\OtherTok}[1]{\textcolor[rgb]{0.56,0.35,0.01}{#1}}
\newcommand{\FunctionTok}[1]{\textcolor[rgb]{0.00,0.00,0.00}{#1}}
\newcommand{\VariableTok}[1]{\textcolor[rgb]{0.00,0.00,0.00}{#1}}
\newcommand{\ControlFlowTok}[1]{\textcolor[rgb]{0.13,0.29,0.53}{\textbf{#1}}}
\newcommand{\OperatorTok}[1]{\textcolor[rgb]{0.81,0.36,0.00}{\textbf{#1}}}
\newcommand{\BuiltInTok}[1]{#1}
\newcommand{\ExtensionTok}[1]{#1}
\newcommand{\PreprocessorTok}[1]{\textcolor[rgb]{0.56,0.35,0.01}{\textit{#1}}}
\newcommand{\AttributeTok}[1]{\textcolor[rgb]{0.77,0.63,0.00}{#1}}
\newcommand{\RegionMarkerTok}[1]{#1}
\newcommand{\InformationTok}[1]{\textcolor[rgb]{0.56,0.35,0.01}{\textbf{\textit{#1}}}}
\newcommand{\WarningTok}[1]{\textcolor[rgb]{0.56,0.35,0.01}{\textbf{\textit{#1}}}}
\newcommand{\AlertTok}[1]{\textcolor[rgb]{0.94,0.16,0.16}{#1}}
\newcommand{\ErrorTok}[1]{\textcolor[rgb]{0.64,0.00,0.00}{\textbf{#1}}}
\newcommand{\NormalTok}[1]{#1}
\usepackage{graphicx,grffile}
\makeatletter
\def\maxwidth{\ifdim\Gin@nat@width>\linewidth\linewidth\else\Gin@nat@width\fi}
\def\maxheight{\ifdim\Gin@nat@height>\textheight\textheight\else\Gin@nat@height\fi}
\makeatother
% Scale images if necessary, so that they will not overflow the page
% margins by default, and it is still possible to overwrite the defaults
% using explicit options in \includegraphics[width, height, ...]{}
\setkeys{Gin}{width=\maxwidth,height=\maxheight,keepaspectratio}
\IfFileExists{parskip.sty}{%
\usepackage{parskip}
}{% else
\setlength{\parindent}{0pt}
\setlength{\parskip}{6pt plus 2pt minus 1pt}
}
\setlength{\emergencystretch}{3em}  % prevent overfull lines
\providecommand{\tightlist}{%
  \setlength{\itemsep}{0pt}\setlength{\parskip}{0pt}}
\setcounter{secnumdepth}{5}
% Redefines (sub)paragraphs to behave more like sections
\ifx\paragraph\undefined\else
\let\oldparagraph\paragraph
\renewcommand{\paragraph}[1]{\oldparagraph{#1}\mbox{}}
\fi
\ifx\subparagraph\undefined\else
\let\oldsubparagraph\subparagraph
\renewcommand{\subparagraph}[1]{\oldsubparagraph{#1}\mbox{}}
\fi

% set default figure placement to htbp
\makeatletter
\def\fps@figure{htbp}
\makeatother

\usepackage{lscape}
\usepackage{longtable}

\title{Overview of the \texttt{NetSim} package}
\author{Gustavo A. Ballen\footnote{Museu de Zoologia da Universidade de São
  Paulo}}
\date{}

\begin{document}
\maketitle

{
\hypersetup{linkcolor=black}
\setcounter{tocdepth}{3}
\tableofcontents
}
\begin{Shaded}
\begin{Highlighting}[]
\KeywordTok{library}\NormalTok{(NetSim)  }\CommentTok{# simulate bd time-tree networks: https://github.com/jjustison/NetSim}
\KeywordTok{library}\NormalTok{(TreeSim)  }\CommentTok{# simulate bd time-trees}
\end{Highlighting}
\end{Shaded}

\begin{verbatim}
Loading required package: ape
\end{verbatim}

\begin{verbatim}
Loading required package: geiger
\end{verbatim}

\begin{verbatim}
Registered S3 method overwritten by 'geiger':
  method            from
  unique.multiPhylo ape 
\end{verbatim}

\begin{Shaded}
\begin{Highlighting}[]
\KeywordTok{library}\NormalTok{(ape)  }\CommentTok{# tools for handling phylo objects}
\KeywordTok{library}\NormalTok{(diversitree)  }\CommentTok{# tools for estimating lambda and mu}
\KeywordTok{library}\NormalTok{(lattice)  }\CommentTok{# stacked plots}
\KeywordTok{set.seed}\NormalTok{(}\DecValTok{1000}\NormalTok{)}
\NormalTok{numbsim <-}\StringTok{ }\DecValTok{100}
\end{Highlighting}
\end{Shaded}

\section{The NetSim package}\label{the-netsim-package}

The package builds on the approach of Stadler for \texttt{TreeSim}, but
with modifications by Zhang et al. (2018) for including hybridization as
part of the birth-death (BD) stochastic process. The functions are named
similarly to those in \texttt{TreeSim} but with \texttt{bdh} insted of
just \texttt{bd}.

The package allows trees conditional to time or number of species, as
well as the typical sampling and extinction patterns allowed by
\texttt{TreeSim}. However, it allows to model hybridization in three
important ways:

\begin{itemize}
\tightlist
\item
  Effect of the hybridization on generation of lineages (positive,
  negative, neutral)
\item
  Hybridization probability as a function of genetic distance (the
  closer two lineages the easier for them to hybridize)
\item
  Hybridization probability covariating with a trait along the phylogeny
  (for example, if hybridization is dependent on the two parents having
  a similar trait)
\end{itemize}

Simulations are controlled via two main parameters: Origin time, and
number of species. For the former, a given time interval is specified
for the simulation to run, leaving \(n\) as a free quantity. On the
other hand, the SSE and GSE approaches fix the final number of species
\(n\) and allow the simulation to run an unconstrained amount of time
until such quantity is attained.

\section{General network simulation
features}\label{general-network-simulation-features}

Comparing the manual entry for the functions \texttt{sim.bdh.age},
\texttt{sim.bdh.taxa.ssa}, and \texttt{sim.bdh.taxa.gsa}, we find that
they differ in the very first arguments (Table \ref{tab:sim_args}).
These first arguments (\texttt{age}, \texttt{m}, and \texttt{n}) control
the endpoint of the simulation. All of the remaining arguments specify
the model oh birth-death hybridization. All functions return an object
of either class \texttt{phylo} or \texttt{multiphylo} depending on how
many trees were requested in the argument \texttt{numbsim}.

Below we will explore each of these functions and their appropriate
context.

\section{Preliminary conditions: Yule and BD
trees}\label{preliminary-conditions-yule-and-bd-trees}

Let's see if the functions produce reliable trees under Yule and BD
without hybridization. Parameter estimation on the tree samples is
inspired in this blog post by
\href{http://www.phytools.org/Cordoba2017/ex/12/BD-models.html}{Revell
and Harmon}

\subsection{Yule trees}\label{yule-trees}

First let's simulate the same tree samples using \texttt{TreeSim} as a
control for comparing the results by \texttt{NetSim}:

\begin{Shaded}
\begin{Highlighting}[]
\NormalTok{yuleTS <-}\StringTok{ }\KeywordTok{sim.bd.age}\NormalTok{(}\DataTypeTok{age =} \DecValTok{4}\NormalTok{, }\DataTypeTok{numbsim =}\NormalTok{ numbsim, }\DataTypeTok{lambda =} \DecValTok{1}\NormalTok{, }\DataTypeTok{mu =} \DecValTok{0}\NormalTok{, }\DataTypeTok{frac =} \DecValTok{1}\NormalTok{, }\DataTypeTok{mrca =} \OtherTok{FALSE}\NormalTok{, }
    \DataTypeTok{complete =} \OtherTok{TRUE}\NormalTok{, }\DataTypeTok{K =} \DecValTok{0}\NormalTok{)}

\NormalTok{## remove any trees with no taxa}
\NormalTok{yuleTS <-}\StringTok{ }\NormalTok{yuleTS[}\OperatorTok{!}\KeywordTok{sapply}\NormalTok{(}\DataTypeTok{X =}\NormalTok{ yuleTS, }\DataTypeTok{FUN =}\NormalTok{ is.null)]}
\NormalTok{## remove non-phylo elements}
\NormalTok{yuleTS <-}\StringTok{ }\NormalTok{yuleTS[}\KeywordTok{sapply}\NormalTok{(}\DataTypeTok{X =}\NormalTok{ yuleTS, }\DataTypeTok{FUN =}\NormalTok{ is.phylo)]}
\NormalTok{## remove any trees with <10 taxa}
\NormalTok{yuleTS <-}\StringTok{ }\NormalTok{yuleTS[}\KeywordTok{sapply}\NormalTok{(}\DataTypeTok{X =}\NormalTok{ yuleTS, }\DataTypeTok{FUN =} \ControlFlowTok{function}\NormalTok{(x) }\KeywordTok{length}\NormalTok{(x}\OperatorTok{$}\NormalTok{tip.label) }\OperatorTok{>=}\StringTok{ }\DecValTok{10}\NormalTok{)]}
\CommentTok{# fit a Yule model with diversitree}
\NormalTok{yuleTSFits <-}\StringTok{ }\KeywordTok{lapply}\NormalTok{(}\DataTypeTok{X =}\NormalTok{ yuleTS, }\DataTypeTok{FUN =} \ControlFlowTok{function}\NormalTok{(x) diversitree}\OperatorTok{::}\KeywordTok{make.yule}\NormalTok{(}\DataTypeTok{tree =}\NormalTok{ x, }
    \DataTypeTok{sampling.f =} \DecValTok{1}\NormalTok{))}
\CommentTok{# fetch the coefficients estimates on the previous lik funs}
\NormalTok{yuleTSMLEcoefs <-}\StringTok{ }\KeywordTok{lapply}\NormalTok{(}\DataTypeTok{X =}\NormalTok{ yuleTSFits, }\DataTypeTok{FUN =} \ControlFlowTok{function}\NormalTok{(x) }\KeywordTok{coef}\NormalTok{(diversitree}\OperatorTok{::}\KeywordTok{find.mle}\NormalTok{(}\DataTypeTok{func =}\NormalTok{ x, }
    \DataTypeTok{x.init =} \FloatTok{0.5}\NormalTok{)))}
\end{Highlighting}
\end{Shaded}

Now we will simulate a Yule tree sample using \texttt{NetSim}:

\begin{Shaded}
\begin{Highlighting}[]
\CommentTok{# age = 4 should generate moderate-sized trees mu = 0 no extinction nu = 0 no}
\CommentTok{# hybridization hybprobs = c(0.5, 0.5, 0.5) no prior difference among types of}
\CommentTok{# effects of hyb on descendant lineages hyb.inher.fxn beta(1,1) -> uniform(0,1),}
\CommentTok{# no pattern of inheritance, frac = 1 we'll simulate complete genealogies}
\CommentTok{# complete = TRUE we don't have extinction, this plays no role hyb.rate.fxn =}
\CommentTok{# NULL we don't want hybridization to be a function of genetic distance}
\CommentTok{# trait.model = NULL we don't need a covariate to include in this simple model}
\NormalTok{yuleNS <-}\StringTok{ }\NormalTok{NetSim}\OperatorTok{::}\KeywordTok{sim.bdh.age}\NormalTok{(}\DataTypeTok{age =} \DecValTok{4}\NormalTok{, }\DataTypeTok{numbsim =}\NormalTok{ numbsim, }\DataTypeTok{lambda =} \DecValTok{1}\NormalTok{, }\DataTypeTok{mu =} \DecValTok{0}\NormalTok{, }\DataTypeTok{nu =} \DecValTok{0}\NormalTok{, }
    \DataTypeTok{hybprops =} \KeywordTok{c}\NormalTok{(}\FloatTok{0.5}\NormalTok{, }\FloatTok{0.5}\NormalTok{, }\FloatTok{0.5}\NormalTok{), }\DataTypeTok{hyb.inher.fxn =}\NormalTok{ NetSim}\OperatorTok{::}\KeywordTok{make.beta.draw}\NormalTok{(}\DecValTok{1}\NormalTok{, }\DecValTok{1}\NormalTok{), }\DataTypeTok{frac =} \DecValTok{1}\NormalTok{, }
    \DataTypeTok{mrca =} \OtherTok{TRUE}\NormalTok{, }\DataTypeTok{complete =} \OtherTok{TRUE}\NormalTok{, }\DataTypeTok{stochsampling =} \OtherTok{FALSE}\NormalTok{, }\DataTypeTok{hyb.rate.fxn =} \OtherTok{NULL}\NormalTok{, }\DataTypeTok{trait.model =} \OtherTok{NULL}\NormalTok{)}

\NormalTok{## remove any trees with no taxa}
\NormalTok{yuleNS <-}\StringTok{ }\NormalTok{yuleNS[}\OperatorTok{!}\KeywordTok{sapply}\NormalTok{(}\DataTypeTok{X =}\NormalTok{ yuleNS, }\DataTypeTok{FUN =}\NormalTok{ is.null)]}
\NormalTok{## remove non-phylo elements}
\NormalTok{yuleNS <-}\StringTok{ }\NormalTok{yuleNS[}\KeywordTok{sapply}\NormalTok{(}\DataTypeTok{X =}\NormalTok{ yuleNS, }\DataTypeTok{FUN =}\NormalTok{ is.phylo)]}
\NormalTok{## remove any trees with <10 taxa}
\NormalTok{yuleNS <-}\StringTok{ }\NormalTok{yuleNS[}\KeywordTok{sapply}\NormalTok{(}\DataTypeTok{X =}\NormalTok{ yuleNS, }\DataTypeTok{FUN =} \ControlFlowTok{function}\NormalTok{(x) }\KeywordTok{length}\NormalTok{(x}\OperatorTok{$}\NormalTok{tip.label) }\OperatorTok{>=}\StringTok{ }\DecValTok{10}\NormalTok{)]}
\CommentTok{# fit a Yule model with diversitree}
\NormalTok{yuleNSFits <-}\StringTok{ }\KeywordTok{lapply}\NormalTok{(}\DataTypeTok{X =}\NormalTok{ yuleNS, }\DataTypeTok{FUN =} \ControlFlowTok{function}\NormalTok{(x) diversitree}\OperatorTok{::}\KeywordTok{make.yule}\NormalTok{(}\DataTypeTok{tree =}\NormalTok{ x, }
    \DataTypeTok{sampling.f =} \DecValTok{1}\NormalTok{))}
\CommentTok{# fetch the coefficients estimates on the previous lik funs}
\NormalTok{yuleNSMLEcoefs <-}\StringTok{ }\KeywordTok{lapply}\NormalTok{(}\DataTypeTok{X =}\NormalTok{ yuleNSFits, }\DataTypeTok{FUN =} \ControlFlowTok{function}\NormalTok{(x) }\KeywordTok{coef}\NormalTok{(diversitree}\OperatorTok{::}\KeywordTok{find.mle}\NormalTok{(}\DataTypeTok{func =}\NormalTok{ x, }
    \DataTypeTok{x.init =} \FloatTok{0.5}\NormalTok{)))}
\end{Highlighting}
\end{Shaded}

Plot both distributions for visual comparison. Please note that the
original value of \(\lambda\) used in the simulations is represented by
the vertical black line:

\begin{Shaded}
\begin{Highlighting}[]
\NormalTok{yuleLambda <-}\StringTok{ }\KeywordTok{data.frame}\NormalTok{(}\DataTypeTok{lambda =} \KeywordTok{c}\NormalTok{(}\KeywordTok{unlist}\NormalTok{(yuleTSMLEcoefs), }\KeywordTok{unlist}\NormalTok{(yuleNSMLEcoefs)), }
    \DataTypeTok{method =} \KeywordTok{c}\NormalTok{(}\KeywordTok{rep}\NormalTok{(}\StringTok{"TreeSim"}\NormalTok{, }\DataTypeTok{times =} \KeywordTok{length}\NormalTok{(yuleTSMLEcoefs)), }\KeywordTok{rep}\NormalTok{(}\StringTok{"NetSim"}\NormalTok{, }\DataTypeTok{times =} \KeywordTok{length}\NormalTok{(yuleNSMLEcoefs))))}
\NormalTok{yuleLambdaHist <-}\StringTok{ }\NormalTok{lattice}\OperatorTok{::}\KeywordTok{densityplot}\NormalTok{(}\OperatorTok{~}\NormalTok{lambda }\OperatorTok{|}\StringTok{ }\NormalTok{method, }\DataTypeTok{data =}\NormalTok{ yuleLambda, }\DataTypeTok{layout =} \KeywordTok{c}\NormalTok{(}\DecValTok{1}\NormalTok{, }
    \DecValTok{2}\NormalTok{), }\DataTypeTok{plot.points =} \OtherTok{FALSE}\NormalTok{, }\DataTypeTok{panel =} \ControlFlowTok{function}\NormalTok{(x, y, ...) \{}
    \KeywordTok{panel.densityplot}\NormalTok{(x, ...)}
    \KeywordTok{panel.abline}\NormalTok{(}\DataTypeTok{v =} \DecValTok{1}\NormalTok{)}
\NormalTok{\})}
\NormalTok{my_theme <-}\StringTok{ }\KeywordTok{trellis.par.get}\NormalTok{()}
\NormalTok{my_theme}\OperatorTok{$}\NormalTok{strip.background}\OperatorTok{$}\NormalTok{col <-}\StringTok{ "grey80"}
\NormalTok{my_theme}\OperatorTok{$}\NormalTok{plot.symbol}\OperatorTok{$}\NormalTok{col <-}\StringTok{ "grey60"}
\NormalTok{my_theme}\OperatorTok{$}\NormalTok{plot.polygon}\OperatorTok{$}\NormalTok{col <-}\StringTok{ "grey90"}

\NormalTok{yuleLHplot <-}\StringTok{ }\KeywordTok{update}\NormalTok{(yuleLambdaHist, }\DataTypeTok{par.settings =}\NormalTok{ my_theme)}
\KeywordTok{print}\NormalTok{(yuleLHplot)}
\end{Highlighting}
\end{Shaded}

\includegraphics{test_netsim_files/figure-latex/unnamed-chunk-4-1.pdf}

\subsection{Birth-death trees}\label{birth-death-trees}

\begin{Shaded}
\begin{Highlighting}[]
\CommentTok{# mrca = FALSE because the tree needs to be binary. This param behaves the}
\CommentTok{# opposite in NetSim!  complete = FALSE drops extinct tips b/c diversitree needs}
\CommentTok{# an ultrametric tree}
\NormalTok{bdTS <-}\StringTok{ }\KeywordTok{sim.bd.age}\NormalTok{(}\DataTypeTok{age =} \FloatTok{9.3}\NormalTok{, }\DataTypeTok{numbsim =}\NormalTok{ numbsim, }\DataTypeTok{lambda =} \FloatTok{0.9}\NormalTok{, }\DataTypeTok{mu =} \FloatTok{0.5}\NormalTok{, }\DataTypeTok{frac =} \FloatTok{0.9}\NormalTok{, }
    \DataTypeTok{mrca =} \OtherTok{FALSE}\NormalTok{, }\DataTypeTok{complete =} \OtherTok{FALSE}\NormalTok{, }\DataTypeTok{K =} \DecValTok{0}\NormalTok{)}

\NormalTok{## remove any trees with no taxa}
\NormalTok{bdTS <-}\StringTok{ }\NormalTok{bdTS[}\OperatorTok{!}\KeywordTok{sapply}\NormalTok{(}\DataTypeTok{X =}\NormalTok{ bdTS, }\DataTypeTok{FUN =}\NormalTok{ is.null)]}
\NormalTok{## remove non-phylo elements}
\NormalTok{bdTS <-}\StringTok{ }\NormalTok{bdTS[}\KeywordTok{sapply}\NormalTok{(}\DataTypeTok{X =}\NormalTok{ bdTS, }\DataTypeTok{FUN =}\NormalTok{ is.phylo)]}
\NormalTok{## remove any trees with <10 taxa}
\NormalTok{bdTS <-}\StringTok{ }\NormalTok{bdTS[}\KeywordTok{sapply}\NormalTok{(}\DataTypeTok{X =}\NormalTok{ bdTS, }\DataTypeTok{FUN =} \ControlFlowTok{function}\NormalTok{(x) }\KeywordTok{length}\NormalTok{(x}\OperatorTok{$}\NormalTok{tip.label) }\OperatorTok{>=}\StringTok{ }\DecValTok{10}\NormalTok{)]}
\CommentTok{# fit a Bd model with diversitree}
\NormalTok{bdTSFits <-}\StringTok{ }\KeywordTok{lapply}\NormalTok{(}\DataTypeTok{X =}\NormalTok{ bdTS, }\DataTypeTok{FUN =} \ControlFlowTok{function}\NormalTok{(x) diversitree}\OperatorTok{::}\KeywordTok{make.bd}\NormalTok{(}\DataTypeTok{tree =}\NormalTok{ x, }\DataTypeTok{sampling.f =} \FloatTok{0.6}\NormalTok{))}
\CommentTok{# fetch the coefficients estimates on the previous lik funs}
\NormalTok{bdTSMLEcoefs <-}\StringTok{ }\KeywordTok{lapply}\NormalTok{(}\DataTypeTok{X =}\NormalTok{ bdTSFits, }\DataTypeTok{FUN =} \ControlFlowTok{function}\NormalTok{(x) }\KeywordTok{coef}\NormalTok{(diversitree}\OperatorTok{::}\KeywordTok{find.mle}\NormalTok{(}\DataTypeTok{func =}\NormalTok{ x, }
    \DataTypeTok{x.init =} \KeywordTok{c}\NormalTok{(}\FloatTok{0.5}\NormalTok{, }\FloatTok{0.5}\NormalTok{))))}
\end{Highlighting}
\end{Shaded}

\begin{verbatim}
Warning in mle.search(func2, x.init, control, lower, upper): Convergence
problems in find.mle: code = 3 (see ?nlm for details)
\end{verbatim}

Now we will simulate a BD tree sample using \texttt{NetSim}:

\begin{Shaded}
\begin{Highlighting}[]
\CommentTok{# age = 9 should generate moderate-sized trees nu = 0 no hybridization hybprobs =}
\CommentTok{# c(0.5, 0.5, 0.5) no prior difference among types of effects of hyb on}
\CommentTok{# descendant lineages hyb.inher.fxn beta(1,1) -> uniform(0,1), no pattern of}
\CommentTok{# inheritance, frac = 1 we'll simulate complete genealogies mrca = TRUE because}
\CommentTok{# the tree needs to be binary complete = FALSE drop extinct tips b/c diversitree}
\CommentTok{# needs an ultrametric tree hyb.rate.fxn = NULL we don't want hybridization to be}
\CommentTok{# a function of genetic distance trait.model = NULL we don't need a covariate to}
\CommentTok{# include in this simple model}
\NormalTok{bdNS <-}\StringTok{ }\NormalTok{NetSim}\OperatorTok{::}\KeywordTok{sim.bdh.age}\NormalTok{(}\DataTypeTok{age =} \FloatTok{9.3}\NormalTok{, }\DataTypeTok{numbsim =}\NormalTok{ numbsim, }\DataTypeTok{lambda =} \FloatTok{0.9}\NormalTok{, }\DataTypeTok{mu =} \FloatTok{0.5}\NormalTok{, }
    \DataTypeTok{nu =} \DecValTok{0}\NormalTok{, }\DataTypeTok{hybprops =} \KeywordTok{c}\NormalTok{(}\FloatTok{0.5}\NormalTok{, }\FloatTok{0.5}\NormalTok{, }\FloatTok{0.5}\NormalTok{), }\DataTypeTok{hyb.inher.fxn =}\NormalTok{ NetSim}\OperatorTok{::}\KeywordTok{make.beta.draw}\NormalTok{(}\DecValTok{1}\NormalTok{, }
        \DecValTok{1}\NormalTok{), }\DataTypeTok{frac =} \FloatTok{0.9}\NormalTok{, }\DataTypeTok{mrca =} \OtherTok{TRUE}\NormalTok{, }\DataTypeTok{complete =} \OtherTok{FALSE}\NormalTok{, }\DataTypeTok{stochsampling =} \OtherTok{FALSE}\NormalTok{, }\DataTypeTok{hyb.rate.fxn =} \OtherTok{NULL}\NormalTok{, }
    \DataTypeTok{trait.model =} \OtherTok{NULL}\NormalTok{)}

\NormalTok{## remove any trees with no taxa}
\NormalTok{bdNS <-}\StringTok{ }\NormalTok{bdNS[}\OperatorTok{!}\KeywordTok{sapply}\NormalTok{(}\DataTypeTok{X =}\NormalTok{ bdNS, }\DataTypeTok{FUN =}\NormalTok{ is.null)]}
\NormalTok{## remove non-phylo elements}
\NormalTok{bdNS <-}\StringTok{ }\NormalTok{bdNS[}\KeywordTok{sapply}\NormalTok{(}\DataTypeTok{X =}\NormalTok{ bdNS, }\DataTypeTok{FUN =}\NormalTok{ is.phylo)]}
\NormalTok{## remove any trees with <10 taxa}
\NormalTok{bdNS <-}\StringTok{ }\NormalTok{bdNS[}\KeywordTok{sapply}\NormalTok{(}\DataTypeTok{X =}\NormalTok{ bdNS, }\DataTypeTok{FUN =} \ControlFlowTok{function}\NormalTok{(x) }\KeywordTok{length}\NormalTok{(x}\OperatorTok{$}\NormalTok{tip.label) }\OperatorTok{>=}\StringTok{ }\DecValTok{10}\NormalTok{)]}
\CommentTok{# remove non-binary}
\NormalTok{bdNS <-}\StringTok{ }\NormalTok{bdNS[}\KeywordTok{sapply}\NormalTok{(}\DataTypeTok{X =}\NormalTok{ bdNS, }\DataTypeTok{FUN =}\NormalTok{ is.binary)]}
\CommentTok{# fit a Bd model with diversitree}
\NormalTok{bdNSFits <-}\StringTok{ }\KeywordTok{lapply}\NormalTok{(}\DataTypeTok{X =}\NormalTok{ bdNS, }\DataTypeTok{FUN =} \ControlFlowTok{function}\NormalTok{(x) diversitree}\OperatorTok{::}\KeywordTok{make.bd}\NormalTok{(}\DataTypeTok{tree =}\NormalTok{ x, }\DataTypeTok{sampling.f =} \FloatTok{0.6}\NormalTok{))}
\CommentTok{# fetch the coefficients estimates on the previous lik funs}
\NormalTok{bdNSMLEcoefs <-}\StringTok{ }\KeywordTok{lapply}\NormalTok{(}\DataTypeTok{X =}\NormalTok{ bdNSFits, }\DataTypeTok{FUN =} \ControlFlowTok{function}\NormalTok{(x) }\KeywordTok{coef}\NormalTok{(diversitree}\OperatorTok{::}\KeywordTok{find.mle}\NormalTok{(}\DataTypeTok{func =}\NormalTok{ x, }
    \DataTypeTok{x.init =} \KeywordTok{c}\NormalTok{(}\FloatTok{0.5}\NormalTok{, }\FloatTok{0.5}\NormalTok{))))}
\end{Highlighting}
\end{Shaded}

Plot both distributions for visual comparisonPlease note that the
original value of \(\lambda\) and \(\mu\) used in the simulations is
represented by the vertical black line:

\begin{Shaded}
\begin{Highlighting}[]
\NormalTok{bdParams <-}\StringTok{ }\KeywordTok{data.frame}\NormalTok{(}\DataTypeTok{lambda =} \KeywordTok{c}\NormalTok{(}\KeywordTok{sapply}\NormalTok{(bdTSMLEcoefs, }\ControlFlowTok{function}\NormalTok{(x) }\KeywordTok{unlist}\NormalTok{(x)[}\StringTok{"lambda"}\NormalTok{]), }
    \KeywordTok{sapply}\NormalTok{(bdNSMLEcoefs, }\ControlFlowTok{function}\NormalTok{(x) }\KeywordTok{unlist}\NormalTok{(x)[}\StringTok{"lambda"}\NormalTok{])), }\DataTypeTok{mu =} \KeywordTok{c}\NormalTok{(}\KeywordTok{sapply}\NormalTok{(bdTSMLEcoefs, }
    \ControlFlowTok{function}\NormalTok{(x) }\KeywordTok{unlist}\NormalTok{(x)[}\StringTok{"mu"}\NormalTok{]), }\KeywordTok{sapply}\NormalTok{(bdNSMLEcoefs, }\ControlFlowTok{function}\NormalTok{(x) }\KeywordTok{unlist}\NormalTok{(x)[}\StringTok{"mu"}\NormalTok{])), }
    \DataTypeTok{method =} \KeywordTok{c}\NormalTok{(}\KeywordTok{rep}\NormalTok{(}\StringTok{"TreeSim"}\NormalTok{, }\DataTypeTok{times =} \KeywordTok{length}\NormalTok{(bdTSMLEcoefs)), }\KeywordTok{rep}\NormalTok{(}\StringTok{"NetSim"}\NormalTok{, }\DataTypeTok{times =} \KeywordTok{length}\NormalTok{(bdNSMLEcoefs))))}

\CommentTok{# build plot}
\NormalTok{bdLambdaHist <-}\StringTok{ }\NormalTok{lattice}\OperatorTok{::}\KeywordTok{densityplot}\NormalTok{(}\OperatorTok{~}\NormalTok{lambda }\OperatorTok{|}\StringTok{ }\NormalTok{method, }\DataTypeTok{data =}\NormalTok{ bdParams, }\DataTypeTok{layout =} \KeywordTok{c}\NormalTok{(}\DecValTok{1}\NormalTok{, }
    \DecValTok{2}\NormalTok{), }\DataTypeTok{plot.points =} \OtherTok{FALSE}\NormalTok{, }\DataTypeTok{panel =} \ControlFlowTok{function}\NormalTok{(x, y, ...) \{}
    \KeywordTok{panel.densityplot}\NormalTok{(x, ...)}
    \KeywordTok{panel.abline}\NormalTok{(}\DataTypeTok{v =} \FloatTok{0.9}\NormalTok{)}
\NormalTok{\})}
\NormalTok{my_theme <-}\StringTok{ }\KeywordTok{trellis.par.get}\NormalTok{()}
\NormalTok{my_theme}\OperatorTok{$}\NormalTok{strip.background}\OperatorTok{$}\NormalTok{col <-}\StringTok{ "grey80"}
\NormalTok{my_theme}\OperatorTok{$}\NormalTok{plot.symbol}\OperatorTok{$}\NormalTok{col <-}\StringTok{ "grey60"}
\NormalTok{my_theme}\OperatorTok{$}\NormalTok{plot.polygon}\OperatorTok{$}\NormalTok{col <-}\StringTok{ "grey90"}

\NormalTok{bdLHplot <-}\StringTok{ }\KeywordTok{update}\NormalTok{(bdLambdaHist, }\DataTypeTok{par.settings =}\NormalTok{ my_theme)}
\KeywordTok{print}\NormalTok{(bdLHplot)}
\end{Highlighting}
\end{Shaded}

\includegraphics{test_netsim_files/figure-latex/unnamed-chunk-7-1.pdf}

\begin{Shaded}
\begin{Highlighting}[]
\NormalTok{bdMuHist <-}\StringTok{ }\NormalTok{lattice}\OperatorTok{::}\KeywordTok{densityplot}\NormalTok{(}\OperatorTok{~}\NormalTok{mu }\OperatorTok{|}\StringTok{ }\NormalTok{method, }\DataTypeTok{data =}\NormalTok{ bdParams, }\DataTypeTok{layout =} \KeywordTok{c}\NormalTok{(}\DecValTok{1}\NormalTok{, }\DecValTok{2}\NormalTok{), }
    \DataTypeTok{plot.points =} \OtherTok{FALSE}\NormalTok{, }\DataTypeTok{panel =} \ControlFlowTok{function}\NormalTok{(x, y, ...) \{}
        \KeywordTok{panel.densityplot}\NormalTok{(x, ...)}
        \KeywordTok{panel.abline}\NormalTok{(}\DataTypeTok{v =} \FloatTok{0.5}\NormalTok{)}
\NormalTok{    \})}
\NormalTok{my_theme <-}\StringTok{ }\KeywordTok{trellis.par.get}\NormalTok{()}
\NormalTok{my_theme}\OperatorTok{$}\NormalTok{strip.background}\OperatorTok{$}\NormalTok{col <-}\StringTok{ "grey80"}
\NormalTok{my_theme}\OperatorTok{$}\NormalTok{plot.symbol}\OperatorTok{$}\NormalTok{col <-}\StringTok{ "grey60"}
\NormalTok{my_theme}\OperatorTok{$}\NormalTok{plot.polygon}\OperatorTok{$}\NormalTok{col <-}\StringTok{ "grey90"}

\NormalTok{bdMHplot <-}\StringTok{ }\KeywordTok{update}\NormalTok{(bdMuHist, }\DataTypeTok{par.settings =}\NormalTok{ my_theme)}
\KeywordTok{print}\NormalTok{(bdMHplot)}
\end{Highlighting}
\end{Shaded}

\includegraphics{test_netsim_files/figure-latex/unnamed-chunk-7-2.pdf}

\subsection{Remarks}\label{remarks}

The results for Yule trees are equivalent for both methods. However,
estimation of \(\lambda\) and \(\mu\) are 1) not completely equivalent
as judged from the plots, and 2) more seriously, do not match the
starting values in the simulations, that is, \texttt{lambda\ =\ 0.9} and
\texttt{mu\ =\ 0.5}.

This might be a limitation of the method used by \texttt{diversitree}
for fitting birth-death models (said to be Nee et al., 1994), but if so
we need a better way to assess with back-comparison how good are these
two estimators and whether they are in fact equivalent.

Once sorted this out we can consider to use NetSim for networks in the
specific case of time-trees, or more specifically, birth-death
time-trees.

\section{\texorpdfstring{Simulation spanning a fixed time interval:
\texttt{sim.bdh.age}}{Simulation spanning a fixed time interval: sim.bdh.age}}\label{simulation-spanning-a-fixed-time-interval-sim.bdh.age}

The plot engine \texttt{NetSim::plottable.net} is said to reorient the
tree for improved plotting with respect to hybrid edges. However, this
also creates extinct tips for hybridization events (see the plot on the
right below, sister to \texttt{t11} and compare with the same node in
the plot to the left).

This does not happen all the time, and this is the best combination of
parameters that allows to reconstruct a reproducible example of such
behavior. Why do this happen? Does it represent a difference in the
nature of the hybridization event? Note that some other hybrid edges do
not induce an ``extinct'' node sister to the hybrid edge.

\begin{Shaded}
\begin{Highlighting}[]
\NormalTok{age_fixed <-}\StringTok{ }\KeywordTok{sim.bdh.age}\NormalTok{(}\DataTypeTok{age =} \FloatTok{2.3}\NormalTok{, }\DataTypeTok{numbsim =} \DecValTok{1}\NormalTok{, }\DataTypeTok{lambda =} \FloatTok{0.45}\NormalTok{, }\DataTypeTok{mu =} \DecValTok{0}\NormalTok{, }\DataTypeTok{nu =} \FloatTok{0.15}\NormalTok{, }
    \DataTypeTok{hybprops =} \KeywordTok{c}\NormalTok{(}\FloatTok{0.5}\NormalTok{, }\FloatTok{0.5}\NormalTok{, }\FloatTok{0.5}\NormalTok{), }\DataTypeTok{hyb.inher.fxn =} \KeywordTok{make.beta.draw}\NormalTok{(}\DecValTok{1}\NormalTok{, }\DecValTok{1}\NormalTok{), }\DataTypeTok{frac =} \DecValTok{1}\NormalTok{, }
    \DataTypeTok{mrca =} \OtherTok{TRUE}\NormalTok{, }\DataTypeTok{complete =} \OtherTok{TRUE}\NormalTok{, }\DataTypeTok{stochsampling =} \OtherTok{FALSE}\NormalTok{, }\DataTypeTok{hyb.rate.fxn =} \OtherTok{NULL}\NormalTok{, }\DataTypeTok{trait.model =} \OtherTok{NULL}\NormalTok{)}
\KeywordTok{str}\NormalTok{(age_fixed)}
\end{Highlighting}
\end{Shaded}

\begin{verbatim}
List of 1
 $ :List of 6
  ..$ edge        : num [1:38, 1:2] 13 13 15 15 16 16 17 17 20 20 ...
  ..$ tip.label   : chr [1:12] "t1" "t16" "t2" "t7" ...
  ..$ edge.length : num [1:38] 1.1498 0.0905 0.4748 0.4955 0.7861 ...
  ..$ Nnode       : num 27
  ..$ reticulation: num [1:8, 1:2] 21 23 18 23 29 33 25 33 26 27 ...
  .. ..- attr(*, "dimnames")=List of 2
  .. .. ..$ : NULL
  .. .. ..$ : chr [1:2] "from" "to"
  ..$ inheritance : num [1:8] 0.183 0.127 0.876 0.782 0.472 ...
  ..- attr(*, "class")= chr [1:2] "evonet" "phylo"
\end{verbatim}

\begin{Shaded}
\begin{Highlighting}[]
\KeywordTok{par}\NormalTok{(}\DataTypeTok{mfrow =} \KeywordTok{c}\NormalTok{(}\DecValTok{1}\NormalTok{, }\DecValTok{2}\NormalTok{))}
\KeywordTok{plot}\NormalTok{(age_fixed[[}\DecValTok{1}\NormalTok{]])}
\end{Highlighting}
\end{Shaded}

\begin{verbatim}
Warning in min(x): no non-missing arguments to min; returning Inf
\end{verbatim}

\begin{verbatim}
Warning in max(x): no non-missing arguments to max; returning -Inf
\end{verbatim}

\begin{verbatim}
Warning in min(x): no non-missing arguments to min; returning Inf
\end{verbatim}

\begin{verbatim}
Warning in max(x): no non-missing arguments to max; returning -Inf
\end{verbatim}

\begin{Shaded}
\begin{Highlighting}[]
\KeywordTok{plot}\NormalTok{(}\KeywordTok{plottable.net}\NormalTok{(age_fixed[[}\DecValTok{1}\NormalTok{]]))}
\end{Highlighting}
\end{Shaded}

\includegraphics{test_netsim_files/figure-latex/unnamed-chunk-8-1.pdf}

\begin{small}
\begin{landscape}
\begin{longtable}[c]{p{2cm}p{1.2cm}p{1.2cm}p{1cm}p{2.6cm}p{2.6cm}p{2.6cm}p{5cm}}
\caption{Summary of arguments for simulation functions in \texttt{NetSim}. X Represents arguments present but without a default value. Filled cells otherwise represent the default values in each function, and empty cells represent arguments absent from a given function.}
\label{tab:sim_args}\\
\multicolumn{1}{c}{\textbf{argument}} & \textbf{arg.type} & \textbf{arg.str} & \textbf{arg.len} & \textbf{sim.bdh.age} & \textbf{sim.bdh.taxa.ssa} & \textbf{sim.bdh.taxa.gsa} & \multicolumn{1}{c}{\textbf{notes}} \\
\endfirsthead
%
\multicolumn{8}{c}%
{{\bfseries Table \thetable\ continued from previous page}} \\
\multicolumn{1}{c}{\textbf{argument}} & \textbf{arg.type} & \textbf{arg.str} & \textbf{arg.len} & \textbf{sim.bdh.age} & \textbf{sim.bdh.taxa.ssa} & \textbf{sim.bdh.taxa.gsa} & \multicolumn{1}{c}{\textbf{notes}} \\
\endhead
%
\hline
\texttt{age} & double & vector & 1 & X &  &  & Total simulation time length \\
\texttt{m} & integer & vector & 1 &  &  & X & Number of taxa to simulate under SSA \\
\texttt{n} & integer & vector & 1 &  & X & X & Number of final taxa \\
\texttt{numbsim} & double & vector & 1 & X & X & X & Number of networks to simulate \\
\texttt{lambda} & double & vector & 1 & X & X & X & Speciation rate \\
\texttt{mu} & double & vector & 1 & X & X & X & Extinction rate \\
\texttt{nu} & double & vector & 1 & X & X & X & Hybridization rate \\
\texttt{hybprops} & double & vector & 3 & X & X & X & Proportion of hybridizations that are lineage generative, lineage degenerative, and lineage neutral respectively \\
\texttt{hyb.inher.fxn} & closure & function & 1 & X & X & X & Function for drawing hybrid inheritance probabilities . Uses a function that draws from \texttt{rDIST}-type functions, e.g., \texttt{make.beta.draw} \\
\texttt{frac} & double & vector & 1 & \texttt{1} & \texttt{1} & \texttt{1} & Proportion of species sampled included in the phylogeny \\
\texttt{mrca} & boolean & vector & 1 & \texttt{FALSE} & \texttt{FALSE} & \texttt{FALSE} & Either origin time counts from the initial point (one lineage, \texttt{FALSE}) or from the first split (two lineages, \texttt{TRUE}) \\
\texttt{complete} & boolean & vector & 1 & \texttt{TRUE} & \texttt{TRUE} & \texttt{TRUE} & Either a tree preserving extinct lineages (if \texttt{TRUE}) or droping such extinct lineages (if \texttt{FALSE}) \\
\texttt{stochsampling} & boolean & vector & 1 & \texttt{FALSE} & \texttt{FALSE} & \texttt{FALSE} & If \texttt{TRUE}, each tip is included into the final tree with probability defined by the argument \texttt{frac}. Otherwise every lineage is guaranteed to reach the present because the phylogeny is defined as complete if \texttt{frac = 1.0} \\
\texttt{hyb.rate.fxn} & closure & function & 1 & \texttt{NULL} & \texttt{NULL} & \texttt{NULL} & Uses a function that draws from density-type functions, e.g., \texttt{make.exp.decay} \\
\texttt{trait.model} & multiple & list & 5 & \texttt{NULL} & \texttt{NULL} & \texttt{NULL} & The same structure as \texttt{make.polyploid.model}
\end{longtable}
\end{landscape}
\end{small}

\bibliography{/home/balleng/Dropbox/Gustavo/Bibtex/PhDThesis.bib}

\end{document}
